\section{Question 6: Which Tax Regime Does the Household Prefer?}

\subsection*{Method}

To determine which tax system the household prefers, we compare the beginning-of-life value function $V_0(K_{1,0}, K_{2,0})$ under joint and individual taxation. The value function gives the expected discounted lifetime utility when the household starts at a given level of human capital and behaves optimally from that point on:
%
\begin{equation}
    V_0(K_{1,0}, K_{2,0}) = \max_{\{h_{1,t}, h_{2,t}\}_{t=0}^{T-1}} \sum_{t=0}^{T-1} \beta^t \, U(c_t, h_{1,t}, h_{2,t})
\end{equation}
%
Here $U$ is the household objective evaluated at $\mu = 0.5$. A higher $V_0$ means the household is better off. We use $\lambda^{\text{indiv},\star} = 1.7670$ for the individual tax system so that the comparison is revenue-neutral, meaning any difference in welfare comes from the tax structure itself, not from a change in the overall tax burden.

We get $V_0$ by solving each model and then using bilinear interpolation of the solution array at the relevant initial state.

Comparing raw value-function levels reveals the \emph{sign} of the welfare difference but not its size, because utils have no natural scale. Following the compensating-variation logic of the course welfare analysis, we therefore also report how much extra initial human capital member~1 would need under \emph{joint} taxation to be as well off as under (revenue-neutral) individual taxation:
%
\begin{equation}
    \Delta K_1^{\star}(K_{1,0},K_{2,0}) = \bigl\{\, x : V_0^{\text{joint}}(K_{1,0}+x,\,K_{2,0}) = V_0^{\text{indiv}}(K_{1,0},K_{2,0}) \,\bigr\}.
\end{equation}
%
This is the simple-model counterpart of the asset-compensation measure $a^{\star}$ from the lecture --- here the household has no savings, so the relevant endowment is human capital --- and it translates the util gap into an equivalent first-period wage gain of $\exp(\alpha_{1,1}\,\Delta K_1^{\star})-1$.

\subsection*{Case 1: Asymmetric Couple ($K_{1,0} = 2$, $K_{2,0} = 0$)}

\begin{center}
\begin{tabular}{l r}
\hline
$V_0$ under joint taxation      & $-0.2267$ \\
$V_0$ under individual taxation & $-0.2267$ \\
\hline
\textbf{Preferred regime}       & Indifferent \\
\hline
\end{tabular}
\end{center}

For the asymmetric couple the two value functions agree to four decimals, and the compensating variation is only $\Delta K_1^{\star} = 0.009$ units of human capital --- equivalent to raising member~1's wage by about $0.09\%$. The household is therefore effectively \textbf{indifferent} between the regimes. Intuitively, with $K_{2,0}=0$ the secondary earner barely works, so the channel through which individual taxation helps --- a lower marginal rate on the second income --- is almost entirely shut off, and a revenue-neutral reform leaves welfare essentially untouched.

\subsection*{Case 2: Symmetric Couple ($K_{1,0} = 8$, $K_{2,0} = 8$)}

\begin{center}
\begin{tabular}{l r}
\hline
$V_0$ under joint taxation      & $-0.1891$ \\
$V_0$ under individual taxation & $-0.1883$ \\
\hline
\textbf{Preferred regime}       & Individual \\
\hline
\end{tabular}
\end{center}

For the symmetric high-capital couple, the preference for individual taxation is clearer. Both members have high earnings, so under joint taxation their combined income pushes them into higher marginal tax brackets. Under individual taxation, each is taxed on their own income separately, which results in lower marginal rates and higher welfare. The gain is now economically meaningful: the compensating variation is $\Delta K_1^{\star} = 0.30$ units of human capital, equivalent to about a $3.0\%$ higher wage for member~1.

\subsection*{Discussion}

In both cases the household (weakly) prefers individual taxation. The intuition is that joint taxation creates a penalty for two-earner households because pooling incomes pushes the household into higher tax brackets. Individual taxation removes this penalty. The effect is stronger for the symmetric couple where both members earn a lot, because the ``marriage tax'' from income pooling is larger when both incomes are high.

These conclusions are robust to the discretisation: refining the human-capital grid from $N_k = 20$ to $N_k = 30$ leaves the preference ranking unchanged and the compensating variation for the symmetric couple stable at $\Delta K_1^{\star} \approx 0.30$, confirming that the welfare differences are economic rather than numerical artefacts.
